\begin{table}[t!]
\caption{What each breadth encoder responds to. Mean cosine distance between a context and
a minimally edited variant of it, over 3,000 target contexts.}
\label{tab:probe}
\small
\begin{tabular}{@{}lccc@{}}
\toprule
Encoder & Target swapped & Control word swapped & Ratio \\
\midrule
XL-LEXEME & $.1197$ & $.0009$ & $135$$\times$ \\
MPNet & $.1313$ & $.0106$ & $12$$\times$ \\
\bottomrule
\end{tabular}
\par\vspace{3pt}
\parbox{\linewidth}{\footnotesize Note. The target swap replaces the target term with the
other concept's term; the control swap replaces an unrelated content word in the same
context. Both conditions use the same contexts. A higher ratio means the representation is
more specific to the target and less contaminated by the rest of the sentence. Across the
11 annual series the two encoders' breadth estimates correlate at
mean $r = .20$.}
\end{table}
